\begin{tablalafrox}{Diccionario del Producto 1: Gestión del proyecto. Fuente: según cita por cuenta en cada encabezado.}{tab:T14-dic-prod-1}
  {L{1.2cm} Y Y Y C{0.9cm} L{2.3cm}}
  {\cab{Código} & \cab{Paquete de trabajo} & \cab{Entregable} & \cab{Criterio de aceptación} & \cab{Resp.} & \cab{Período}}
  \multicolumn{6}{>{\columncolor{lafroxBg}\RaggedRight\arraybackslash}p{\dimexpr\textwidth-2\tabcolsep\relax}}{\textbf{Cuenta 1.1 — Constitución y confirmaciones contractuales.} Esta cuenta de control reúne los tres documentos con que el proyecto arranca formalmente. {\color{lafroxGris}Fuente: (Bases Administrativas, arts. 17.2–17.6 y Formulario T-14, p. 63); (Bases Técnicas Transversales, secciones 17.2–17.6, pp. 32–34).}} \\
  1.1.1 & Acta de constitución del proyecto, firmada con la Contraparte Técnica & Acta firmada por LafroX y Puelche. & Está firmada en el mes 1 e identifica al patrocinador, a la Contraparte Técnica, al Jefe de Proyecto y los objetivos. & JP & Ambas · Mes 1. \\
  1.1.2 & Documento de validación con el CLIENTE de los problemas y metas del proyecto & Documento de validación con los cambios acordados, si los hay. & Las gerencias del CLIENTE confirman por escrito los problemas y las metas. & JP & Ambas · Mes 1. \\
  1.1.3 & Acta de confirmación con el CLIENTE de la fecha de inicio, la evidencia para el proveedor de lácteos, las ventanas de paso a producción y el acuse tributario & Acta con la respuesta del CLIENTE a cada tema. & Los cuatro temas tienen una respuesta escrita incorporada al plan. Para el proveedor de lácteos, el CLIENTE valida el formato y el canal de un expediente inicial: registros disponibles de recepción y lote, recorrido hasta despacho y lista de destinatarios recuperables, conciliados con los respaldos existentes. El acta registra vacíos y responsables de completarlos. Hasta habilitar 5.2.1, Calidad mantiene la evidencia con los sistemas y registros existentes. La aceptación del proveedor se documenta por separado. & JP & Ambas · Mes 1, antes del H1. \\
  \midrule
  \multicolumn{6}{>{\columncolor{lafroxBg}\RaggedRight\arraybackslash}p{\dimexpr\textwidth-2\tabcolsep\relax}}{\textbf{Cuenta 1.2 — Planificación y líneas base.} Esta cuenta de control reúne los documentos que dicen cómo y cuándo se hará el trabajo, y los informes que muestran si se cumple. {\color{lafroxGris}Fuente: (Bases Administrativas, arts. 17.2–17.6 y Formulario T-14, p. 63); (Bases Técnicas Transversales, secciones 17.2–17.6, pp. 32–34).}} \\
  1.2.1 & Plan para la dirección del proyecto, con el calendario de congelamientos de septiembre, diciembre y el cierre mensual & Plan para la dirección con el calendario de congelamientos. & Ningún paso a producción, corte de datos ni despliegue cae en una fecha prohibida; los meses 16 y 21 coinciden con el Art. 17°; el plan está aprobado. & JP & Ambas · Mes 1 (H1). Se actualiza con cada cambio. \\
  1.2.2 & EDT y diccionario de la EDT (línea base del alcance) & EDT dibujada y diccionario aprobados. & Todo requerimiento de la matriz (2.1.3) cae en un paquete, y ningún paquete queda sin responsable. & JP & Ambas · Mes 2 (H1). \\
  1.2.3 & Carta Gantt de 56 meses con los hitos H1 a H12 y la ruta crítica & Carta Gantt aprobada. & Cubre los 56 meses, respeta el Art. 17° y nombra la ruta crítica con su cálculo. & JP & Ambas · Mes 2. \\
  1.2.4 & Documento de nivelación de recursos y frentes del solapamiento de los meses 13 a 15 y 19 a 20 & Nivelación de recursos por etapa y por frente. & Ningún mes asigna a una persona por sobre su capacidad, y los dos frentes del solapamiento tienen dotación propia. & JP & Ambas · Mes 2. Se revisa antes del mes 13. \\
  1.2.5 & Reunión semanal de seguimiento con su informe semanal & Una reunión y un informe por semana. & La reunión se hace cada semana, y su informe se publica con el avance por paquete y las desviaciones frente al cronograma. & JP & Ambas · Semanal, durante toda la implementación. \\
  \midrule
  \multicolumn{6}{>{\columncolor{lafroxBg}\RaggedRight\arraybackslash}p{\dimexpr\textwidth-2\tabcolsep\relax}}{\textbf{Cuenta 1.3 — Interesados y comunicaciones.} Esta cuenta de control reúne el trabajo con las personas que pueden apoyar o frenar el proyecto, y el lugar donde se toman las decisiones. {\color{lafroxGris}Fuente: (Bases Administrativas, arts. 17.2–17.6 y Formulario T-14, p. 63); (Bases Técnicas Transversales, secciones 17.2–17.6, pp. 32–34).}} \\
  1.3.1 & Registro de interesados con su influencia, interés y estrategia & Registro de interesados. & Están los 17 grupos, cada uno con un responsable de la relación, y el registro se revisa antes de cada ola de implantación. & JP & Ambas · Mes 1. \\
  1.3.2 & Plan de comunicaciones con calendario y canales por grupo de interés & Plan de comunicaciones. & Cada grupo del registro tiene su mensaje, su canal y su frecuencia. & JP & Ambas · Mes 1. \\
  1.3.3 & Registro de solicitudes de cambio con análisis de impacto & Registro de solicitudes de cambio con su análisis de impacto. & Cada solicitud tiene su análisis de impacto y la decisión del Comité Ejecutivo, y ningún cambio se ejecuta sin aprobación. & JP & Ambas · Durante todo el contrato. \\
  \midrule
  \multicolumn{6}{>{\columncolor{lafroxBg}\RaggedRight\arraybackslash}p{\dimexpr\textwidth-2\tabcolsep\relax}}{\textbf{Cuenta 1.4 — Gestión de la calidad.} Esta cuenta de control reúne las reglas que impiden que algo defectuoso llegue a los usuarios. {\color{lafroxGris}Fuente: (Bases Administrativas, arts. 17.2–17.6 y Formulario T-14, p. 63); (Bases Técnicas Transversales, secciones 17.2–17.6, pp. 32–34).}} \\
  1.4.1 & Plan de calidad y plan de pruebas & Plan de calidad y plan de pruebas aprobados. & Cada tipo de prueba que exigen las Bases tiene su ambiente, su calendario y su criterio de salida. & CAL & Ambas · Mes 2. \\
  1.4.2 & Puertas de calidad automáticas: cobertura de lógica de negocio del 70 \%, global del 80 \% y del código modificado del 80 \% & Puertas activas con informe separado para cada ámbito. & Se bloquea una versión si incumple cualquiera de los tres mínimos: lógica de negocio 70 \% (RT-04.11), global 80 \% y código modificado 80 \% (políticas de calidad de la oferta). Los denominadores y resultados se informan por separado; ninguna versión llega a Preproducción con pruebas en falla. & CAL & Ambas · Desde que existe la cadena de integración (3.1.4). \\
  \midrule
  \multicolumn{6}{>{\columncolor{lafroxBg}\RaggedRight\arraybackslash}p{\dimexpr\textwidth-2\tabcolsep\relax}}{\textbf{Cuenta 1.5 — Gestión de riesgos.} Esta cuenta de control mantiene vivo el registro de riesgos y verifica que cada respuesta se ejecute. {\color{lafroxGris}Fuente: (Bases Administrativas, arts. 17.2–17.6 y Formulario T-14, p. 63); (Bases Técnicas Transversales, secciones 17.2–17.6, pp. 32–34).}} \\
  1.5.1 & Plan de riesgos y registro inicial de riesgos según ISO 31000 & Plan de riesgos y registro inicial. & Cada riesgo tiene dueño, disparador y respuesta, y cada respuesta tiene un paquete o una actividad en la EDT. & JP & Ambas · Mes 2. \\
  1.5.2 & Registro de riesgos revisado en cada Comité de Proyecto & Registro actualizado con su historial. & El registro se revisa en cada Comité de Proyecto, y cada cambio de evaluación queda registrado. & JP & Ambas · Quincenal, durante toda la implementación. \\
  \midrule
  \multicolumn{6}{>{\columncolor{lafroxBg}\RaggedRight\arraybackslash}p{\dimexpr\textwidth-2\tabcolsep\relax}}{\textbf{Cuenta 1.6 — Plan de Reversibilidad.} Esta cuenta reúne el Plan de Reversibilidad y su actualización anual para que Puelche pueda cambiar de proveedor o asumir la operación sin interrumpir el servicio (BA, art. 77, p. 39). {\color{lafroxGris}Fuente: (Bases Administrativas, arts. 17.2–17.6 y Formulario T-14, p. 63); (Bases Técnicas Transversales, secciones 17.2–17.6, pp. 32–34).}} \\
  1.6.1 & Plan de Reversibilidad inicial (primeros 90 días) & Plan de Reversibilidad. & Se entrega dentro de los primeros 90 días con sus cuatro contenidos. & JP & Ambas · Hasta el mes 3. \\
  1.6.2 & Primera actualización anual del Plan de Reversibilidad durante la implementación & Plan actualizado. & El inventario coincide con lo instalado a la fecha de la actualización. & JP & Ambas · Mes 15, al cumplirse un año de la entrega de 1.6.1; las actualizaciones de los meses 27, 39 y 51 son el paquete 1.6.3 (Art. 77.2). \\
  1.6.3 & Actualizaciones anuales del Plan de Reversibilidad durante los 36 meses de operación & Tres actualizaciones del plan, una por año de operación. & En cada actualización, el inventario, las claves y la exportación están vigentes. & JP & Operación · Meses 27, 39 y 51, continuando la serie anual que inicia 1.6.2 (Art. 77.2). \\
  \midrule
  \multicolumn{6}{>{\columncolor{lafroxBg}\RaggedRight\arraybackslash}p{\dimexpr\textwidth-2\tabcolsep\relax}}{\textbf{Cuenta 1.7 — Gobierno, control y cierre.} Esta cuenta reúne las instancias de gobierno obligatorias para ambas partes (BA, art. 71, p. 37), los mecanismos de control y reporte (BTT, cap. 19) y el cierre de la implementación. Las actas de todos los comités se levantan dentro de los dos días hábiles siguientes, con acuerdos, responsables y plazos (RT-19.09). {\color{lafroxGris}Fuente: (Bases Administrativas, arts. 17.2–17.6 y Formulario T-14, p. 63); (Bases Técnicas Transversales, secciones 17.2–17.6, pp. 32–34).}} \\
  1.7.1 & Oficina de gestión del proyecto con metodología basada en el PMBOK & Oficina constituida, con su manual de procedimientos y sus plantillas. & Los procesos de alcance, cronograma, costos, riesgos, cambios y comunicaciones tienen procedimiento y responsable desde el mes 1. & JP & Ambas · Mes 1, y durante toda la implementación. \\
  1.7.2 & Actas del Comité Ejecutivo mensual, con la aprobación de los cambios de alcance & Un acta por reunión mensual, con las decisiones y los cambios aprobados o rechazados. & Hay un acta por mes, levantada dentro de dos días hábiles, y todo cambio ejecutado tiene su aprobación previa en un acta. & JP & Ambas · Mensual, durante todo el contrato. \\
  1.7.3 & Actas del Comité de Proyecto quincenal & Un acta por reunión quincenal. & Hay un acta cada quince días, levantada dentro de dos días hábiles. & JP & Ambas · Quincenal, durante la implementación. \\
  1.7.4 & Actas del Comité de Arquitectura mensual & Un acta por reunión mensual, con las decisiones de arquitectura registradas. & Hay un acta por mes, y cada decisión aprobada está en el registro de decisiones de arquitectura. & ARQ & Ambas · Mensual, durante todo el contrato. \\
  1.7.5 & Actas del Comité de Operación mensual, desde el mes 13 & Un acta por reunión mensual, desde el mes 13 hasta el mes 56. & Hay un acta por mes, levantada dentro de dos días hábiles. & SRE & Ambas · Mensual, meses 13 a 56. \\
  1.7.6 & Informe mensual de avance con valor ganado & Un informe por mes. & Cada informe trae los índices de valor ganado, y cada desviación mayor al 10 \% tiene su aviso y su plan dentro de cinco días hábiles. & JP & Ambas · Mensual, durante la implementación. \\
  1.7.7 & Tablero de estado del proyecto disponible para el CLIENTE & Tablero publicado, con acceso para el CLIENTE. & El CLIENTE accede en todo momento, y los datos coinciden con el último informe. & JP & Ambas · Desde el mes 2. \\
  1.7.8 & Espacio colaborativo con la documentación, los entregables, las actas, los riesgos y los cambios & Espacio colaborativo habilitado, con acceso para el CLIENTE. & Cada entregable aceptado y cada acta están en el espacio, en su versión vigente. & JP & Ambas · Desde el mes 1. \\
  1.7.9 & Informe de lecciones aprendidas y archivo de la implementación & Informe de lecciones aprendidas y archivo. & El informe cubre las dos etapas, y el archivo contiene todas las actas de H1 a H12. & JP & Ambas · Después del H12. \\
  \midrule
  \multicolumn{6}{>{\columncolor{lafroxBg}\RaggedRight\arraybackslash}p{\dimexpr\textwidth-2\tabcolsep\relax}}{\textbf{Cuenta 1.8 — Cumplimiento normativo y contractual.} Esta cuenta de control reúne las obligaciones legales y contractuales que no producen software, pero que deben mantenerse y acreditarse: cumplimiento normativo, datos personales, certificación de seguridad, garantías, seguros y obligaciones laborales. {\color{lafroxGris}Fuente: (Bases Administrativas, arts. 17.2–17.6 y Formulario T-14, p. 63); (Bases Técnicas Transversales, secciones 17.2–17.6, pp. 32–34).}} \\
  1.8.1 & Matriz de cumplimiento normativo por obligación legal, con su control y su evidencia & Matriz de cumplimiento normativo. & Cada obligación aplicable tiene su control y su evidencia, y la matriz está aprobada por el CLIENTE. & SEG & Ambas · Mes 1, y actualizada ante cada cambio normativo. \\
  1.8.2 & Acuerdo de tratamiento de datos personales con el CLIENTE & Acuerdo firmado. & Está firmado antes de tratar datos reales; es decir, antes de cargar el maestro (5.4.2). & SEG & Ambas · Mes 1. \\
  1.8.3 & Registro de actividades de tratamiento y evaluación de impacto en protección de datos & Registro de tratamientos y evaluaciones de impacto. & Cada tratamiento tiene su registro; los de alto riesgo tienen evaluación de impacto; hay una contraparte designada. & SEG & Ambas · Mes 2, y actualizado ante cada tratamiento nuevo. \\
  1.8.4 & Certificación ISO/IEC 27001 vigente y entrega anual de los informes de certificación & Certificación vigente durante el contrato e informe anual de certificaciones. & La certificación nunca está vencida, y hay un informe entregado cada año. & SEG & Ambas · Anual, durante todo el contrato. \\
  1.8.5 & Garantía de fiel cumplimiento, seguros y certificados laborales vigentes durante el contrato & Garantía constituida, pólizas acreditadas cada año y certificados laborales mensuales. & Ninguna garantía ni póliza queda vencida, y hay un certificado laboral por mes. & JP & Ambas · Desde la adjudicación hasta el término del contrato. \\
\end{tablalafrox}
